\chapter{Quantum Group Representations}
\label{ch:quantum-group-reps}

The two-stage factorisation $\Phi_d = \SpCh_{\Sigma_{d-1}, C} \circ \PhiFA_d$ of Chapter~\ref{ch:e1-chiral} (Theorem~\ref{thm:phi-two-stage-factorisation}) sends a CY$_d$
category $\cC$ to a chiral algebra $\Phi_d(\cC)$ on a reference curve~$C$.
At $d \geq 3$ this output is natively $\Eone$; the $\Etwo$-braided
structure belongs to the Drinfeld centre
$\cZ(\Rep^{\Eone}(\Phi_d(\cC)))$
(Chapter~\ref{ch:drinfeld-center}; Remark~\ref{rem:e2-on-drinfeld-centre}),
not to $\Phi_d(\cC)$ itself. Conjecture~CY-C asks when this braided
category is a quantum-group representation category. The constructed
and conditional criteria separate into three loci:
\begin{enumerate}[label=(\alph*)]
 \item \emph{Kazhdan--Lusztig} at $d = 2$: affine vertex algebra
 $V_k(\frakg)$ versus quantum group $\Uq(\frakg)$ at
 $q = e^{\pi i/(k+h^\vee)}$;
 \item \emph{Yangian RTT} from the ordered collision residue
 $r(z)$ of the Vol~I bar complex;
 \item \emph{BPS / CoHA} at $d = 3$: toric Donaldson--Thomas
 invariants assemble into an associative positive half; the full
 quantum group is reached only after adjoining the negative half,
 Cartan part, Hopf pairing, completion, and Drinfeld double.
\end{enumerate}
Their mutual compatibility is the content of CY-C on the constructed
locus and requires additional comparison data elsewhere.


\section{\texorpdfstring{$\Rep_q(\frakg)$}{Rep-q(frakg)} as a braided monoidal category}
\label{sec:rep-q-braided}

\begin{definition}[Category \texorpdfstring{$\Rep_q(\frakg)$}{Rep-q(frakg)}]
\label{def:rep-q}
Let $\frakg$ be a finite-dimensional simple Lie algebra and
$q \in \C^\times$. The category $\Rep_q(\frakg)$ has:
\begin{enumerate}[label=(\roman*)]
 \item \emph{Objects}: finite-dimensional $\Uq(\frakg)$-modules
 (Definition~\ref{def:qgf-drinfeld-jimbo});
 \item \emph{Morphisms}: $\Uq(\frakg)$-linear maps;
 \item \emph{Monoidal structure}: the tensor product
 $V \otimes W$ with $\Uq(\frakg)$-action via the coproduct
 $\Delta$;
 \item \emph{Braiding}: $\sigma_{V,W} = \tau_{V,W} \circ \cR_{V,W}$
 where $\cR$ is the universal $R$-matrix of
 Remark~\ref{rem:qgf-universal-R-formal} and $\tau$ is the
 transposition $v \otimes w \mapsto w \otimes v$.
\end{enumerate}
\end{definition}

\begin{proposition}[Semisimplicity dichotomy]
\label{prop:semisimplicity-dichotomy}
The structure of $\Rep_q(\frakg)$ depends on whether $q$ is generic
or a root of unity:
\begin{enumerate}[label=(\roman*)]
 \item \emph{Generic $q$} (i.e., $q$ not a root of unity):
 $\Rep_q(\frakg)$ is semisimple, with simple objects $V_q(\lambda)$
 parametrized by dominant integral weights $\lambda \in P^+$.
 The characters are the same as the classical case:
 $\mathrm{ch}(V_q(\lambda)) = \mathrm{ch}(V(\lambda))$
 (the Weyl character formula is $q$-independent).
 \item \emph{Root of unity $q = e^{\pi i/(k+h^\vee)}$} in the
 simply-laced normalization, with
 $k \in \Z_{>0}$: $\Rep_q(\frakg)$ is non-semisimple; the
 \emph{fusion category} $\cC_q(\frakg)$ is the quotient by
 negligible morphisms, with finitely many simple objects
 $\{V_q(\lambda) : \lambda \in P^+_k\}$ parametrized by the
 level-$k$ alcove $P^+_k = \{\lambda \in P^+ :
 \langle \lambda, \theta^\vee \rangle \leq k\}$.
\end{enumerate}
This is the Lusztig--Andersen--Paradowski root-of-unity
semisimplification theorem, in the convention where the lacing number
is absorbed into the definition of~$q$.
\end{proposition}

\begin{example}[\texorpdfstring{$\fsl_2$}{fsl-2} at generic \texorpdfstring{$q$}{q}]
\label{ex:sl2-generic}
For $\frakg = \fsl_2$, the simple modules are
$V_q(n)$ ($n \in \Z_{\geq 0}$) with $\dim V_q(n) = n+1$. The
tensor product decomposition is
$V_q(m) \otimes V_q(n) \simeq \bigoplus_{k=|m-n|}^{m+n} V_q(k)$
(same as classical, with $k \equiv m+n \pmod{2}$). The $R$-matrix
on $V_q(1) \otimes V_q(1)$ is
\[
 \cR = q^{1/2}
 \begin{pmatrix}
 q & 0 & 0 & 0 \\
 0 & 1 & q - q^{-1} & 0 \\
 0 & 0 & 1 & 0 \\
 0 & 0 & 0 & q
 \end{pmatrix}
\]
in the basis $\{e_1 \otimes e_1, e_1 \otimes e_2,
e_2 \otimes e_1, e_2 \otimes e_2\}$. This is the standard
solution of the Yang--Baxter equation for the 6-vertex model.
\end{example}

\begin{example}[\texorpdfstring{$\fsl_2$}{fsl-2} at root of unity]
\label{ex:sl2-root-of-unity}
At $q = e^{\pi i/(k+2)}$, the fusion category $\cC_q(\fsl_2)$ has
$k+1$ simple objects $\{V_q(0), V_q(1), \ldots, V_q(k)\}$ with
fusion rules
\[
 V_q(m) \otimes V_q(n)
 = \bigoplus_{\substack{j = |m-n| \\ j \equiv m+n \,(\mathrm{mod}\,2)}}^{\min(m+n, \, 2k-m-n)}
 V_q(j).
\]
The quantum dimension is
$\dim_q V_q(n) = [n+1]_q = (q^{n+1} - q^{-n-1})/(q - q^{-1})$
and the total quantum dimension squared is
$\mathcal{D}^2 = \sum_{j=0}^{k} (\dim_q V_q(j))^2
= (k+2)/(2\sin^2(\pi/(k+2)))$.
The $S$-matrix entries are
$S_{mn} = \sqrt{2/(k+2)} \, \sin(\pi(m+1)(n+1)/(k+2))$.
This is a modular tensor category: the Reshetikhin--Turaev functor
produces the $SU(2)$ Chern--Simons invariants of 3-manifolds.
\end{example}


\section{The \texorpdfstring{$R$}{R}-matrix as categorical \texorpdfstring{$r(z)$}{r(z)}}
\label{sec:r-matrix-categorical}

The $R$-matrix of $\Uq(\frakg)$ is the categorical incarnation of
the collision residue $r(z)$ from the Volume~I bar complex.

\begin{proposition}[\texorpdfstring{$R$}{R}-matrix from bar degree \texorpdfstring{$(1,1)$}{(1,1)}]
\label{prop:r-matrix-bar}
For the affine vertex algebra $V_k(\frakg)$ at level $k$ with
$q = e^{\pi i/(k+h^\vee)}$, the degree-$(1,1)$ component of the
ordered bar complex $B^{\mathrm{ord}}(V_k(\frakg))$ gives the
classical collision residue whose Etingof--Kazhdan quantization is
the universal $R$-matrix of $\Uq(\frakg)$:
\[
 r(z) = \Res^{\mathrm{coll}}_{0,2}(\Theta_{V_k(\frakg)})
 \;\longleftrightarrow\;
 \cR_q = \lim_{z \to 0}\, \cR(z)
\]
where $r(z) = k\,\Omega_{\mathrm{tr}}/z + O(1)$ is the classical $r$-matrix with
Casimir $\Omega \in \frakg \otimes \frakg$ and level prefix $k$,
and $\cR_q$ is the quantized universal $R$-matrix. The spectral
parameter $z$ is the coordinate difference $z = z_1 - z_2$ of two
insertion points on the curve, a rapidity variable in the
integrable lattice interpretation. The pole order of $r(z)$ at
$z = 0$ is one less than that of the OPE, reflecting the fact that
$r(z)$ records the commutator rather than the product.
\end{proposition}

\begin{remark}[Three \texorpdfstring{$r$}{r}-matrices;
\ClaimStatusDefinitional]
\label{rem:three-r-matrices-qg}
Three $r$-matrix-like objects appear:
\begin{enumerate}[label=(\alph*)]
 \item $r^{\mathrm{coll}}(z)$: the bar-intrinsic collision residue
 (Volume~I);
 \item $r(z) = k\,\Omega_{\mathrm{tr}}/z$: the classical $r$-matrix at level $k$
 (Belavin--Drinfeld; vanishes at $k = 0$);
 \item $\cR_q$: the universal $R$-matrix of $\Uq(\frakg)$
 (Drinfeld--Jimbo).
\end{enumerate}
The passage $r(z) \rightsquigarrow \cR_q$ is the Etingof--Kazhdan
quantization theorem. For even $\Einf$-algebras,
$r^{\mathrm{coll}} = r$; for odd generators or genuinely $\Eone$
algebras they can differ.
\end{remark}


\section{Kazhdan--Lusztig equivalence}
\label{sec:kl-equivalence}

The Kazhdan--Lusztig--Finkelberg equivalence is the positive-level
root-of-unity prototype for CY-C.

\begin{theorem}[Kazhdan--Lusztig--Finkelberg;
\ClaimStatusProvedElsewhere]
\label{thm:kl-equivalence}
For a simple Lie algebra $\frakg$ and a positive integer $k$, with
$q = e^{\pi i/(k+h^\vee)}$ in the simply-laced normalization, there
is an equivalence of modular tensor categories
\[
 \mathcal{MTC}_q(\frakg) \;\simeq\; \mathcal{O}^{\mathrm{int}}_k(\hat{\frakg})
\]
where $\mathcal{MTC}_q(\frakg)$ is the semisimplified tilting-module
category of $\Uq(\frakg)$ at a root of unity, and
$\mathcal{O}^{\mathrm{int}}_k(\hat{\frakg})$ is the integrable
highest-weight category of affine Kac--Moody modules at level~$k$
with KZ braiding. In non-simply-laced type the lacing number enters
the denominator of~$q$.
\end{theorem}

\begin{proof}[Attribution]
Kazhdan--Lusztig, \emph{J.\ Amer.\ Math.\ Soc.}~1993--1994, construct
the tensor equivalence for affine Lie algebra categories. Finkelberg,
\emph{GAFA}~1996, proves the positive-level fusion equivalence with the
semisimplified root-of-unity quantum group category. Andersen--Paradowski
give the tilting-module semisimplification used for
$\mathcal{MTC}_q(\frakg)$.
\end{proof}

\begin{remark}[Bar-complex interpretation]
\label{rem:kl-bar-complex}
From the bar-complex perspective, the equivalence states that
the $\Eone$-representation category of the affine vertex algebra
$V_k(\frakg)$ recovers the quantum group representation category.
The Volume~I bar complex $B(V_k(\frakg))$ encodes the quantum group
$R$-matrix in its degree-$(1,1)$ component; the DK bridge (Volume~II,
MC3 on the evaluation-generated core for all simple types) extends this
to the categorical
equivalence.
\end{remark}

\begin{proposition}[KL and the DK bridge]
\label{prop:kl-dk-bridge}
The Kazhdan--Lusztig--Finkelberg equivalence factors through the Drinfeld--Kohno
bridge (Volume~II):
\[
\begin{tikzcd}[column sep=large]
 \mathcal{O}^{\mathrm{int}}_k(\hat{\frakg}) \ar[r, "\sim", "\text{KLF}"'] &
 \mathcal{MTC}_q(\frakg) \\
 \Rep^{\Eone}(V_k(\frakg)) \ar[u, hook] \ar[r, "\text{DK}"] &
 \Rep_q(\frakg) \ar[u, two heads]
\end{tikzcd}
\]
The vertical arrows are: inclusion of integrable affine modules into all
$V_k(\frakg)$-modules (left), and quotient by negligible morphisms
(right). The DK bridge (MC3) provides the horizontal equivalence
at the level of compact objects.
\end{proposition}


\section{Conjecture CY-C: quantum group realization}
\label{sec:conj-cy-c}

\begin{conjecture}[Quantum group realization: Conjecture CY-C]
\label{conj:qg-realization}
For a simple Lie algebra $\frakg$ and $q = e^{\pi i/(k+h^\vee)}$
with $k$ a positive integer, there exists a CY$_2$ category
$\cC(\frakg, q)$ such that:
\begin{enumerate}[label=(\roman*)]
 \item $\Phi(\cC(\frakg, q))$ is an $\Etwo$-chiral algebra whose
 underlying vertex algebra is $V_k(\frakg)$;
 \item $\Rep^{\Etwo}(\Phi(\cC(\frakg, q))) \simeq \Rep_q(\frakg)$
 as braided monoidal categories;
 \item $\kappa_{\mathrm{ch}}(\Phi(\cC(\frakg, q)))
 = \dim(\frakg) \cdot (k + h^\vee)/(2h^\vee)$, recovering the
 Volume~I modular characteristic of $V_k(\frakg)$ as the chiral
 invariant of the $\Phi$-output.
\end{enumerate}
\end{conjecture}

\begin{remark}[Proof inputs for Conjecture CY-C]
\label{rem:cy-c-proof-inputs}
Item (i) is constructed for $\frakg = \fsl_N$ via the Bridgeland
stability conditions on the CY$_2$ resolution of the $A_{N-1}$
surface singularity. Item (ii) at the $\Eone$ level is the content
of MC3 (proved on the evaluation-generated core for all simple types, Volume~I/II). The $\Etwo$
upgrade requires the Drinfeld center passage
(Chapter~\ref{ch:drinfeld-center}). Item (iii) follows from
Theorem~CY-D (Theorem~\ref{thm:cy-modular-characteristic}) when
the CY category is constructed.
\end{remark}


\section{Yangian and RTT realizations}
\label{sec:yangian-rtt}

\begin{definition}[Yangian]
\label{def:yangian}
The \emph{Yangian} $Y(\frakg)$ is the $\C[\hbar]$-algebra with
generators $\{x^{(r)}_i : i = 1, \ldots, \dim \frakg, \;
r = 0, 1, 2, \ldots\}$ and relations determined by the RTT
presentation: $R_{12}(u-v) \, T_1(u) \, T_2(v)
= T_2(v) \, T_1(u) \, R_{12}(u-v)$ where $T(u) = \sum_r T_r u^{-r}$
is the generating matrix and $R_{12}(u)$ is the Yang $R$-matrix.
\end{definition}

\begin{proposition}[Yangian from collision residue]
\label{prop:yangian-collision}
On the evaluation-generated core, the degree-$(1,1)$ collision residue
of the ordered bar complex determines the RTT presentation of the
Yangian:
\[
 Y(\frakg) = \langle r^{\mathrm{coll}}(z) \rangle_{\mathrm{RTT}},
\]
where the RTT relation is the degree-$(1,1,1)$ bar differential after
Etingof--Kazhdan quantization. This is the evaluation-generated-core
content of Volume~II, MC3: the ordered bar complex
$B^{\mathrm{ord}}(V_k(\frakg))$ produces the Yangian presentation from
the modes of $r(z)=k\,\Omega_{\mathrm{tr}}/z$.
\end{proposition}

\begin{remark}[Three distinct operations]
\label{rem:three-operations}
The following three operations on quantum groups must never be
conflated:
\begin{enumerate}[label=(\alph*)]
 \item \emph{Koszul duality}:
 $V_k(\frakg) \mapsto V_k(\frakg)^!$, the Koszul dual vertex
 algebra involving the dual level
 $k' = -k - 2h^\vee$;
 \item \emph{Feigin--Frenkel involution}:
 $k \mapsto -k - 2h^\vee$ within the same family of affine
 algebras; at the critical level $k = -h^\vee$, this produces
 the Feigin--Frenkel center
 $\mathfrak{z}(\hat{\frakg}_{-h^\vee}) = \mathrm{Fun}(
 \mathrm{Op}_{G^L})$;
 \item \emph{Langlands duality}:
 $\frakg \mapsto \frakg^L$ (root-coroot exchange), producing
 a genuinely different Lie algebra for non-simply-laced types.
\end{enumerate}
For $\fsl_2$ (self-dual, self-Langlands), all three coincide on
representation categories. For $G_2$:
$G_2^L = G_2$ (self-Langlands) but the Koszul dual
$V_k(G_2)^!$ is at level $k' = -k - 8$ (since $h^\vee(G_2) = 4$),
while the FF involution also sends $k \mapsto -k-8$; these
coincide as level maps but produce different algebras.
\end{remark}


\section{BPS algebras and quantum groups}
\label{sec:bps-algebras}

The BPS algebra of a CY$_3$ category provides a physical
construction of quantum groups.

\begin{definition}[BPS algebra]
\label{def:bps-algebra}
For a CY$_3$ category $\cC$ with stability condition $\sigma$,
the \emph{BPS algebra} $\cA^{\BPS}_\sigma(\cC)$ is the
cohomological Hall algebra (CoHA) of $\cC$:
\[
 \cA^{\BPS}_\sigma(\cC)
 = \bigoplus_{\gamma \in K_0(\cC)}
 H^*_{\mathrm{BM}}(\cM_\sigma(\gamma), \, \phi_{\mathrm{tr}W})
\]
where $\cM_\sigma(\gamma)$ is the moduli of $\sigma$-semistable
objects of class $\gamma$, $W$ is the CY potential, and
$\phi_{\mathrm{tr}W}$ is the vanishing cycle sheaf. The
multiplication is the Hall product (from short exact sequences).
\end{definition}

\begin{proposition}[CoHA as quantum group]
\label{prop:coha-quantum-group}
For the tripled Jordan quiver (one vertex, three loops $X,Y,Z$) with
potential $W = \mathrm{tr}(XYZ-XZY)$ and CY$_3$ geometry $\C^3$:
\begin{enumerate}[label=(\roman*)]
 \item The CoHA is isomorphic, as an associative Hall algebra, to the positive half of the affine
 Yangian $Y(\widehat{\fgl}_1)^+$ (Schiffmann--Vasserot,
 Maulik--Okounkov);
 \item The full affine Yangian $Y(\widehat{\fgl}_1)$ is recovered only
 after adjoining the negative half, Cartan part, Hopf pairing, and
 completed Drinfeld double of the CoHA;
 \item The quantum toroidal algebra
 $U_{q,t}(\widehat{\widehat{\fgl}}_1)$ arises from the
 $K$-theoretic refinement (replacing cohomology by $K$-theory).
\end{enumerate}
The CoHA itself is not a vertex algebra and not the full affine
Yangian. The $\cW_{1+\infty}$ vertex algebra appears only after the
Drinfeld-double/centre passage and an evaluation representation on a
vacuum module.
\end{proposition}

\begin{proof}[Attribution]
Schiffmann--Vasserot, \emph{Publ.\ IHES}~118 (2013), Theorem~9.1,
identifies the cohomological Hall algebra of the Hilbert scheme model
with the positive half of the affine Yangian of $\widehat{\fgl}_1$ in
the normalisation used for $\C^3$ DT theory. Kapranov--Vasserot,
\emph{Compos.\ Math.}~154 (2018), Theorem~5.1, gives the $K$-theoretic
upgrade. The passage from a positive half to a full quasitriangular
Hopf algebra is the Drinfeld-double construction and requires the
negative half, Cartan part, pairing, and completion.
\end{proof}

\begin{remark}[Slab as bimodule]
\label{rem:bps-slab}
In the Dimofte framework of Volume~II, the BPS algebra
arises from the 3d holomorphic-topological theory on the slab
$X \times [0,1]$. The slab has \emph{two} boundary components
($X \times \{0\}$ and $X \times \{1\}$), making the algebra of
operators on the slab a bimodule for the two boundary algebras.
When the Hall pairing and completion exist, the Drinfeld double is the
endomorphism algebra of the identity bimodule. This bimodule structure
is essential: a Swiss-cheese
disk has one closed and one open boundary component; the slab has
two open boundary components.
\end{remark}

\begin{example}[Resolved conifold]
\label{ex:resolved-conifold}
For the resolved conifold $\cO(-1) \oplus \cO(-1) \to \bP^1$, the
CY$_3$ category is $D^b(\Coh(\cO(-1)^{\oplus 2}))$ and the BPS
algebra is the positive Hall half whose completed Drinfeld double is
the affine Yangian $Y(\widehat{\fsl}_2)$ on the toric chamber. The
chiral modular characteristic is
$\kappa_{\mathrm{ch}} = 1$
(direct McKay computation; the conifold is not a local surface; see
Corollary~\ref{cor:conifold-non-local-surface}). The wall-crossing
formula of Kontsevich--Soibelman governs the dependence of the
BPS spectrum on the stability condition.
\end{example}


\section{Wall-crossing and Kontsevich--Soibelman}
\label{sec:wall-crossing-ks}

The BPS spectrum of a CY$_3$ category depends on the stability
condition $\sigma \in \Stab(\cC)$. As $\sigma$ crosses a wall of
marginal stability, the BPS states reorganize: bound states form or
decay. The Kontsevich--Soibelman wall-crossing formula governs
this reorganization at the level of the motivic Hall algebra.

\begin{definition}[Wall of marginal stability]
\label{def:wall-marginal}
A \emph{wall of marginal stability} in $\Stab(\cC)$ is a real
codimension-$1$ locus $W_{\gamma_1, \gamma_2}$ where two classes
$\gamma_1, \gamma_2 \in K_0(\cC)$ with $\gamma = \gamma_1 + \gamma_2$
have aligned central charges:
$Z_\sigma(\gamma_1)/Z_\sigma(\gamma_2) \in \R_{>0}$.
On opposite sides of $W_{\gamma_1, \gamma_2}$, a bound state of
class $\gamma$ may exist or not.
\end{definition}

\begin{theorem}[Kontsevich--Soibelman wall-crossing formula]
\label{thm:ks-wall-crossing}
For a CY$_3$ category $\cC$ with stability condition $\sigma$,
define the \emph{BPS automorphism} for each class
$\gamma \in K_0(\cC)$:
\[
 U_\gamma = \exp\!\left(
 \sum_{n=1}^{\infty} \frac{(-1)^{n \dim \cM(\gamma)}}{n^2}\,
 \Omega(\gamma; \sigma) \cdot \hat{e}_{n\gamma}
 \right)
\]
where $\Omega(\gamma; \sigma) \in \Z$ is the (numerical)
Donaldson--Thomas invariant and $\hat{e}_\gamma$ is the
corresponding element of the quantum torus algebra
$\hat{e}_\gamma \hat{e}_{\gamma'} = (-1)^{\langle \gamma, \gamma' \rangle}
\hat{e}_{\gamma + \gamma'}$. The ordered product
\[
 \cA_\sigma = \prod_{\gamma : \arg Z_\sigma(\gamma) = \theta}^{\curvearrowleft}
 U_\gamma
\]
(clockwise ordering by phase $\theta = \arg Z_\sigma(\gamma)$) is
\emph{independent of $\sigma$}: as $\sigma$ crosses a wall, the
individual $\Omega(\gamma; \sigma)$ jump, but the ordered product
$\cA_\sigma$ is invariant.
\end{theorem}

\begin{remark}[Motivic refinement]
\label{rem:ks-motivic}
The numerical DT invariants $\Omega(\gamma; \sigma) \in \Z$ are
shadows of the motivic DT invariants
$[\cM_\sigma(\gamma)]_{\mathrm{vir}} \in K_0(\mathrm{Var})$ in the
Grothendieck ring of varieties. The motivic wall-crossing formula
replaces $\Omega$ by the motivic class, and the quantum torus by
the motivic quantum torus. The identification ``scattering =
shadow'' (Volume~I) holds at the motivic level but fails at
the naive BCH level: the full motivic Hall algebra is required.
\end{remark}

\begin{conjecture}[Wall-crossing and the bar complex]
\label{conj:wc-bar}
The wall-crossing formula has a bar-complex interpretation:
crossing a wall $W_{\gamma_1, \gamma_2}$ in $\Stab(\cC)$ induces
a mutation of the ordered bar complex
$B^{\mathrm{ord}}(A_\cC)$ at the corresponding degree.
Concretely:
\begin{enumerate}[label=(\roman*)]
 \item The BPS automorphism $U_\gamma$ corresponds to the
 degree-$|\gamma|$ component of the MC element
 $\Theta_{A_\cC}^{E_1}$ in the $\Eone$ convolution algebra;
 \item Wall-crossing preserves the full MC element
 $\Theta_{A_\cC}^{E_1}$ (it is an invariant of the CY category
 $\cC$, not of the stability condition);
 \item The individual BPS invariants $\Omega(\gamma; \sigma)$ are
 \emph{projections} of $\Theta_{A_\cC}^{E_1}$ that depend on
 $\sigma$; the MC element itself is wall-crossing invariant.
\end{enumerate}
\end{conjecture}


\section{The chiral modular characteristic \texorpdfstring{$\kappa_{\mathrm{ch}}$}{kappa-ch}}
\label{sec:kappa-cat}

The modular characteristic of the quantum group representation
category connects the categorical datum to the genus expansion of
Volume~I.

\begin{proposition}[\texorpdfstring{$\kappa_{\mathrm{ch}}$}{kappa-ch} on the chiral-algebra side of a quantum-group correspondence]
\label{prop:kappa-cat-quantum-groups}
Assume Conjecture~\ref{conj:qg-realization}. For the affine vertex
algebra $V_k(\frakg)$ at level $k$ with
$q = e^{\pi i/(k + h^\vee)}$, the chiral modular characteristic of
$V_k(\frakg)$ (the chiral-algebra image under $\Phi_2$ of the CY$_2$
category $\cC(\frakg, q)$; Conjecture~\ref{conj:qg-realization}) is:
\[
 \kappa_{\mathrm{ch}}(V_k(\frakg))
 = \dim(\frakg) \cdot \frac{k + h^\vee}{2h^\vee}
\]
(Volume~I landscape census). This is the \emph{chiral} invariant,
distinct from the categorical Euler characteristic
$\kappa_{\mathrm{cat}}(\cC) = \chi(\cO_X)$, which may vanish while
$\kappa_{\mathrm{ch}}$ does not. For the standard families:
\begin{center}
\renewcommand{\arraystretch}{1.3}
\begin{tabular}{lll}
 \toprule
 $\frakg$ & $\kappa_{\mathrm{ch}}(V_k(\frakg))$ & $q$ \\
 \midrule
 $\fsl_2$ & $3(k+2)/4$ & $e^{\pi i/(k+2)}$ \\
 $\fsl_N$ & $\frac{N^2-1}{2N}(k+N)$ & $e^{\pi i/(k+N)}$ \\
 $\mathfrak{so}_{2n}$ & $\frac{n(2n-1)(k+2n-2)}{2(2n-2)}$
 & $e^{\pi i/(k+2n-2)}$ \\
 $E_8$ & $\frac{248(k+30)}{60}$
 & $e^{\pi i/(k+30)}$ \\
 \bottomrule
\end{tabular}
\end{center}
\end{proposition}

\begin{remark}[Four \texorpdfstring{$\kappa_\bullet$}{kappa-.}-invariants]
\label{rem:kappa-cat-vs-others}
Four invariants coexist and must be distinguished. Two are
manifold invariants, independent of the algebraization:
\begin{enumerate}[label=(\roman*)]
 \item $\kappa_{\mathrm{cat}}(X) = \chi(\cO_X)$: the holomorphic
 Euler characteristic, a topological invariant;
 \item $\kappa_{\mathrm{fiber}}$: the rank of the relevant lattice
 (e.g.\ Mukai lattice of $H^*(K3, \Z)$).
\end{enumerate}
Two are algebraization invariants, functions of the chiral algebra
produced by a specific construction:
\begin{enumerate}[resume,label=(\roman*)]
 \item $\kappa_{\mathrm{ch}}(A)$: the chiral modular characteristic
 of a chiral algebra $A$, equal to the genus-$1$ Hodge class
 coefficient (Volume~I, Theorem~D);
 \item $\kappa_{\mathrm{BKM}}$: the Borcherds--Kac--Moody modular
 weight, a number-theoretic invariant attached to a BKM algebra
 constructed from a lattice (Chapter~\ref{chap:toroidal-elliptic}).
\end{enumerate}
For $K3 \times E$:
$\kappa_{\mathrm{cat}}(K3 \times E) = \chi(\cO_{K3}) \cdot \chi(\cO_E)
 = 2 \cdot 0 = 0$ by K\"unneth. The compact Hodge/PhiFA scalar is
$\kappa_{\mathrm{ch}}^{\mathrm{Hodge}}(\Phi(K3 \times E))=0$; the additive number
$2+1=3$ is the separate Mukai--Heisenberg specialisation
$\kappa_{\mathrm{ch}}^{\mathrm{Heis}}(K3\times E)$. The BKM weight is
$\kappa_{\mathrm{BKM}}(\Delta_5) = \left.c_N(0)/2\right|_{N=1} = 5$
from the Gritsenko denominator $\Delta_5$ by the universal Borcherds
weight formula (Theorem~\ref{thm:borcherds-weight-kappa-BKM-universal});
the primitive Igusa square is
$\Phi_{10}^{\mathrm{un}}=\Delta_5^2$ of weight~$10$, not the BKM weight.
The identification $\kappa_{\mathrm{ch}} = \chi(\cO_X)$ holds for
CY$_2$ with $h^{1,0} = 0$ (where Serre duality kills the quantum
correction) but fails at odd $d$ by Serre parity
$\chi(\cO_X) = 0$, and fails at $d \geq 3$ in general: the four
invariants are controlled by four different functorialities.
\end{remark}

\begin{proposition}[Chiral scalar complementarity under Feigin--Frenkel involution]
\label{prop:kappa-cat-complementarity}
Under Feigin--Frenkel Koszul duality
$V_k(\frakg) \mapsto V_{k'}(\frakg)$ with $k' = -k - 2h^\vee$, the
chiral modular characteristics of the two chiral-algebra sides satisfy
\[
 \kappa_{\mathrm{ch}}(V_k(\frakg))
 + \kappa_{\mathrm{ch}}(V_{k'}(\frakg))
 = 0,
\]
where on the CY side $q = e^{\pi i/(k + h^\vee)}$ and
$q' = e^{-\pi i/(k + h^\vee)}$. This is the KM/free-field clause of
the Volume~I scalar complementarity (Volume~I, Theorem~C); it is
\emph{not} a statement about the categorical invariant
$\kappa_{\mathrm{cat}}$ (which is a holomorphic Euler characteristic
and behaves additively on Koszul-dual pairs only when the underlying
category is fixed). The $\rho_A \cdot K$ Virasoro/$\cW_N$ clause is
separate.
\end{proposition}

\begin{proof}
Under $k \mapsto k' = -k - 2h^\vee$ the Sugawara-shifted scalar
transforms as
$\kappa_{\mathrm{ch}}(V_{k'}(\frakg))
 = \dim(\frakg)(k' + h^\vee)/(2h^\vee)
 = -\dim(\frakg)(k + h^\vee)/(2h^\vee)
 = -\kappa_{\mathrm{ch}}(V_k(\frakg))$,
so the sum vanishes.
\end{proof}


\section{Shadow obstruction tower for quantum group categories}
\label{sec:shadow-tower-qg}

The shadow obstruction tower $\Theta_{A_\cC}$
(Volume~I) acquires categorical meaning through the
quantum group lens.

\begin{proposition}[Shadow depth from \texorpdfstring{$R$}{R}-matrix pole structure]
\label{prop:shadow-depth-r-matrix}
For $V_k(\frakg)$ with $R$-matrix
$R(z) = 1 + r(z) + O(r^2)$ where $r(z) = \frac{k\,\Omega}{z}$:
\begin{enumerate}[label=(\roman*)]
 \item The shadow depth $r_{\max} = 3$ (class~L): the collision
 residue $r(z)$ has a single pole, the cubic shadow $C$ is
 nonzero (from the Lie bracket), and the quartic resonance class
 $Q$ vanishes by semisimplicity of $\frakg$;
 \item The averaging map (Volume~I) sends
 $r(z) = \frac{k\,\Omega}{z} \mapsto \kappa_{\mathrm{ch}}(V_k(\frakg))$: the full
 $z$-dependent profile is killed by $\Sigma_2$-coinvariance,
 leaving only the scalar modular characteristic;
 \item The KZ associator $\Phi_{\mathrm{KZ}}$ (the degree-$3$
 component of $\Theta_{A_\cC}^{E_1}$) maps under averaging to
 the cubic shadow $C$.
\end{enumerate}
\end{proposition}

\begin{remark}[From quantum groups to the shadow tower]
\label{rem:qg-to-shadow}
The passage from quantum group data to the shadow obstruction tower
inverts the construction of this chapter: the quantum group
$\Uq(\frakg)$ determines the $R$-matrix, which is the degree-$(1,1)$
component of the $\Eone$ MC element $\Theta^{E_1}$, whose
$\Sigma_n$-coinvariant projections are the shadow tower. The
reconstruction problem (from quantum group to chiral algebra) is
solved on the evaluation-generated core by the DK bridge (MC3 for all
simple types on that core): the categorical
data of $\Rep_q(\frakg)$ uniquely determines the chiral algebra
$V_k(\frakg)$ up to the completion issues of MC4.
\end{remark}


\section{Distinctions between Yangian definitions}
\label{sec:two-yangian-definitions}

The Yangian admits two genuinely distinct definitions on a curve.
The weaker one is an $\Eone$-factorisation algebra with RTT
relations; the stronger one supplements this with a chiral
coproduct, spectral $S$-matrix, and the full modular Maurer--Cartan
tower. Equivalence of the two definitions is an open problem at
the heart of the Vol~II E$_1$-core.

\begin{definition}[\texorpdfstring{$\Eone$}{E1}-chiral Yangian (weaker)]
\label{def:e1-chiral-yangian-vol3}
An \emph{$\Eone$-chiral Yangian} on a smooth algebraic curve $X$ is
an $\Eone$-factorisation algebra $\cY^{\mathrm{ch}}$ on $\Ran(X)$
whose local operator algebra, in the RTT presentation, satisfies
the spectral RLL relation
\[
 R_{12}(u-v)\, T_1(u)\, T_2(v)
 = T_2(v)\, T_1(u)\, R_{12}(u-v)
\]
with spectral parameter identified as the coordinate difference
of ordered insertion points: $u = z_1 - z_2$. The factorisation
structure uses \emph{ordered} configuration spaces
$\Conf^{\mathrm{ord}}_n(X)$, and Verdier duality acts by
$R$-matrix inversion $R(u) \mapsto R(u)^{-1}$. This definition is
the Vol~II/I source-of-truth (see Vol~I
\cite{VolI} \ref{def:e1-chiral-yangian-vol3});
the present paragraph records its Vol~III incarnation as input
to $\Phi$.
\end{definition}

\begin{definition}[Chiral Yangian datum (stronger)]
\label{def:chiral-yangian-datum-vol3}
A \emph{chiral Yangian datum} on $X$ is a quadruple
$(A, S(z), \Delta^{\mathrm{ch}}, \{m_k\})$ consisting of:
\begin{enumerate}[label=(\roman*)]
 \item An $\Eone$-chiral algebra $A$ on $\Ran(X)$ (so in
 particular an $\Eone$-chiral Yangian by
 Definition~\ref{def:e1-chiral-yangian-vol3});
 \item A spectral $S$-matrix $S(z) \in \End(A \otimes A)((z))$
 satisfying the spectral YBE;
 \item A chiral coproduct $\Delta^{\mathrm{ch}} \colon A \to
 A \otimes A$ at depth $z = z_1 - z_2$ compatible with the
 $\Eone$-factorisation structure;
 \item Higher operations $\{m_k\}$ generating the modular
 Maurer--Cartan tower $\Theta_A^{E_1}$ in the convolution
 algebra (Volume~I, MC2--MC4).
\end{enumerate}
The four axioms are: Koszul self-recovery
$\Phi(\Phi(A)) \simeq A$ on the Koszul locus; Verdier by
$R$-inversion; existence of a braided-module category
 $\Rep^{\Eone}(A)$; and convergence of the modular MC tower at all
genera.
\end{definition}

\begin{remark}[Forgetful arrow and the open equivalence]
\label{rem:yangian-forgetful-arrow}
There is a forgetful functor from chiral Yangian data to
$\Eone$-chiral Yangians:
$(A, S, \Delta^{\mathrm{ch}}, \{m_k\}) \mapsto A^{E_1}$. Whether
this functor is equivalently surjective is the open equivalence
problem between the two Yangian definitions: given an $\Eone$-chiral Yangian, can the
chiral coproduct $\Delta^{\mathrm{ch}}$ be reconstructed from the
$\Eone$-factorisation data alone? The main obstruction is that
$\Delta^{\mathrm{ch}}$ probes the depth-two collision residue
$r(z_1 - z_2)$ at $z_1 \neq z_2$, while the $\Eone$-factorisation
structure probes only the depth-one collision residue at the
limit $z_1 \to z_2$. Whether the depth-two data is determined by
the depth-one data is unsolved.
\end{remark}

\begin{remark}[Vol III usage convention]
\label{rem:vol3-yangian-convention}
Throughout this chapter and the rest of Vol~III, the term
\emph{Yangian} without further qualifier refers to the
RTT-presented $\Eone$-chiral Yangian of
Definition~\ref{def:e1-chiral-yangian-vol3}. Its local mode algebra is
the RTT algebra of Definition~\ref{def:yangian}; the curve-level object
adds ordered factorisation on~$\Ran(X)$ and therefore is not merely the
classical algebra with notation changed. The full chiral Yangian datum
is invoked explicitly when needed
(e.g.~for the Drinfeld-double construction or when the chiral
coproduct $\Delta^{\mathrm{ch}}$ enters).
\end{remark}


\section{Four Yangian types must be kept distinct}
\label{sec:four-yangian-types}

Four distinct objects are all called ``Yangian'' in the literature;
conflating any two is a type error. Vol~II catalogues the
four-fold typology; the Vol~III incarnation is recorded here.

\begin{definition}[Four Yangian types]
\label{def:four-yangian-types}
The four distinct objects collectively called ``Yangian'' are:
\begin{enumerate}[label=(\roman*)]
 \item \emph{Classical Yangian} $Y_\hbar(\frakg)$: the
 $\C[\hbar]$-Hopf algebra on a point with
 Drinfeld--Jimbo presentation; lives on $\R$ as
 $\Eone$-topological;
 \item \emph{dg-shifted Yangian} $Y_\hbar^{\mathrm{dg}}(\frakg)$:
 the dg-Hopf algebra on a formal disk $D = \Spec\,\C[[z]]$
 obtained as the Etingof--Kazhdan quantization at the formal
 disk;
 \item \emph{Chiral Yangian} $Y(\frakg)^{\mathrm{ch}}$: the
 $\Eone$-chiral algebra on a smooth curve $X$, factorisation
 algebra on $\Ran(X)$ with ordered configuration spaces
 (Definition~\ref{def:e1-chiral-yangian-vol3});
 \item \emph{Spectral Yangian} $Y^{\mathrm{spec}}(\frakg)$: the
 factorisation algebra on the spectral $\A^1_u$ obtained by
 viewing the spectral parameter $u$ as the geometric
 coordinate.
\end{enumerate}
The four are related by functorial passages: classical
$\to$ dg-shifted via formal completion at $\hbar = 0$;
dg-shifted $\to$ chiral via globalisation along the curve;
chiral $\to$ spectral via Fourier-Mukai along the spectral
parameter. None of these passages is an equivalence in
general.
\end{definition}

\begin{remark}[Type errors and where they occur]
\label{rem:yangian-type-errors}
Conflating any two of the four types produces a type error.
Commonly documented confusions:
\begin{enumerate}[label=(\alph*)]
 \item Classical vs dg-shifted: the classical Yangian has no
 dg-structure, the dg-shifted version does. Statements about
 ``the Yangian's dg-Hopf structure'' refer to the dg-shifted
 type only;
 \item Chiral vs spectral: the chiral Yangian lives on the
 \emph{worldsheet curve} $X$; the spectral Yangian lives on
 the \emph{spectral curve} $\A^1_u$. These are different
 spaces with different geometric meaning. Statements like
 ``the chiral algebra has spectral parameter $z$'' must
 specify which curve $z$ parameterizes.
 (The elliptic Belavin caveat from Vol~I applies: the elliptic,
 trigonometric, and rational degenerations refer to the
 spectral parameter $u$, not the worldsheet $z$.)
\end{enumerate}
\end{remark}

\begin{remark}[Vol III scope: chiral Yangian only]
\label{rem:vol3-yangian-scope}
The Yangian in Vol~III is the chiral Yangian
$Y(\frakg)^{\mathrm{ch}}$ of
Definition~\ref{def:four-yangian-types}(iii) throughout. The
classical and dg-shifted Yangians appear only as inputs to the
globalisation procedure constructing
$Y(\frakg)^{\mathrm{ch}}$ on a curve. The spectral Yangian appears
in the Drinfeld-center diagonalisation
(Chapter~\ref{ch:drinfeld-center}) and the K3 toroidal Yangian
(Chapter~\ref{ch:k3-quantum-toroidal}), where the spectral
parameter $u$ identifies with a coordinate on the K3 elliptic
fibration.
\end{remark}


\section{Mukai-Heisenberg lattice sector at \texorpdfstring{$d = 2$}{d = 2}}
\label{sec:mukai-heisenberg-kl}

For K3 surfaces, the cross-volume identification
$\Phi(D^b(\Coh(K3))) = $ $\Etwo$-Mukai-Heisenberg provides a
quantum-group-side lattice criterion: the rank-$24$ Mukai lattice
$\Lambda_{\mathrm{Muk}} = II_{4,20}$ supplies the charge lattice, and the
$R$-matrix at degree $(1,1)$ is the Mukai pairing. This is a
Heisenberg/lattice sector. It is not the $K3\times E$ BKM object and
not the Mukai self-mirror K3 Yangian branch.

\begin{proposition}[Mukai-Heisenberg sector for the lattice \texorpdfstring{$\Lambda_{\mathrm{Muk}}$}{Lambda-Muk}]
\label{prop:kl-k3-lattice}
Let $S$ be a K3 surface with Mukai lattice
$\Lambda_{\mathrm{Muk}} = U^4 \oplus E_8(-1)^2 \simeq II_{4,20}$.
Then:
\begin{enumerate}[label=(\roman*)]
 \item The B-model output
 $\Phi_2(D^b(\Coh(S))) = \mathrm{MH}_2(\Lambda_{\mathrm{Muk}})$
 is the $\Etwo$-Mukai-Heisenberg algebra of
 Definition~\ref{def:e2-mukai-heisenberg};
 \item Its representation category is the Heisenberg Fock category
 graded by $\Lambda_{\mathrm{Muk}}$ and braided by the Mukai pairing;
 \item Since $\Lambda_{\mathrm{Muk}}$ is even unimodular,
 $\Lambda_{\mathrm{Muk}}^{\vee}/\Lambda_{\mathrm{Muk}}=0$. Any finite discriminant-form
 quotient is therefore the trivial pointed category. The usual
 positive-definite rational lattice-VOA modular-tensor-category theorem
 does not apply directly to the indefinite lattice $II_{4,20}$.
\end{enumerate}
\end{proposition}

\begin{proof}
Item~(i) is Vol~III property (U4) of $\Phi$ at $d = 2$
combined with the Hochschild bridge for B-side input
$D^b(\Coh(K3))$. The Heisenberg Fock modules are labelled by charges
in the Mukai lattice, and their braiding is the exponential of the
Mukai pairing. Unimodularity gives
$\Lambda_{\mathrm{Muk}}^{\vee}/\Lambda_{\mathrm{Muk}}=0$, so the discriminant-form quotient
has one simple object. The finite rational lattice-VOA theorem requires
a positive-definite even lattice; here it supplies only the trivial
discriminant check after passing to the quotient, while the full
Heisenberg sector remains a Fock category over the indefinite charge
lattice.
\end{proof}

\begin{remark}[Cross-side comparison: A-side via HMS]
\label{rem:mukai-heisenberg-a-side}
The A-side identification
$\Phi_2(\Fuk(K3)) = \mathrm{MH}_2(\Lambda_{\mathrm{Muk}})$
is Theorem~\ref{thm:fukaya-k3-mukai-heisenberg} of the
preceding Fukaya chapter. Combined with HMS at K3 (Seidel--Smith,
Sheridan; partial), the two sides give the same chiral algebra on the
constructed locus: both Fukaya and derived-coherent inputs to $\Phi_2$
at $d = 2$ produce the $\Etwo$-Mukai-Heisenberg algebra.
\end{remark}


\section{CY-C on constructed loci}
\label{sec:cy-c-constructed-loci}

The quantum-group reconstruction conjecture CY-C asserts that for
a CY$_d$ category $\cC$ with sufficient symmetry, there exists a
generalised quantum group $G(\cC)$. At $d = 2$ the predicted braided
category is $\Rep^{\Etwo}(\Phi_2(\cC))$; at $d \geq 3$ it is
$\cZ(\Rep^{\Eone}(\Phi_d(\cC)))$, the Drinfeld centre of the native
$\Eone$ representation category. CY-C is \textsc{conjectural} in
general; the cases below delimit the constructed boundary.

\begin{conjecture}[CY-C on constructed and unconstructed loci]
\label{conj:cy-c-constructed-loci}
The CY-C claims split as follows:
\begin{enumerate}[label=(\roman*)]
 \item For $d = 2$ with $\cC = D^b(\Coh(K3))$ or $\cC = \Fuk(K3)$:
 the constructed target is the Mukai-Heisenberg lattice sector
 (Proposition~\ref{prop:kl-k3-lattice}); CY-C is verified only at this
 Heisenberg/lattice level;
 \item For $d = 2$ with $\cC = D^b(\Coh(\text{abelian surface}))$:
 $\kappa_{\mathrm{cat}}(A)=\kappa_{\mathrm{ch}}^{\mathrm{Hodge}}(A)=0$
 and $\kappa_{\mathrm{ch}}^{\mathrm{Heis}}(A)=2$. The vanishing of
 $\kappa_{\mathrm{cat}}$ does not construct a trivial quantum group;
 no nontrivial CY-C target is asserted here;
 \item For $d = 2$ with general CY$_2$ $\cC$ admitting an
 $\fsl_N$-symmetry: $G(\cC)$ is conjectured to be
 $U_q(\fsl_N)$ at root of unity, but the construction is
 incomplete beyond the Bridgeland-stability case for
 $A_{N-1}$-singularities (Conjecture~\ref{conj:qg-realization});
 \item For $d = 3$ with $\cC = \Coh(\C^3)$ or
 $\cC = \Coh(\text{conifold})$: the constructed Hall algebra is the
 positive half $Y^+$ of the relevant affine Yangian
 (Proposition~\ref{prop:coha-quantum-group}); the full Yangian and its
 braided representation category require the completed Drinfeld double,
 including negative half, Cartan, Hopf pairing, and centre compatibility;
 \item For $d = 3$ with general compact CY$_3$ $\cC$:
 $G(\cC)$ is \emph{unconstructed}; CY-C remains
 \textsc{conjectural};
 \item For super-CY$_d$ inputs: the super-Yangian
 $Y(\frakg(m|n))$ identification is \textsc{conjectural} at
 the chiral level
 (Vol~III \cite{VolI},
 Definition~\ref{def:chiral-super-yangian-mn});
 \item For $d \geq 4$: $G(\cC)$ is \emph{undefined} in general,
 since any $\Etwo$-braided comparison must pass through
 $\cZ(\Rep^{\Eone}(\Phi(\cC)))$, the Drinfeld centre
 of the $\Eone$-representation category, and the resulting braided
 category may not be of quantum-group type.
\end{enumerate}
\end{conjecture}

\begin{remark}[Pentagon stratification of the six routes]
\label{rem:cy-c-pentagon}
The six routes to $G(K3 \times E)$ (Kummer, Borcherds, MO stable
envelope, McKay, factorization homology, Costello 5d) are six
distinct constructions; only one applies $\Phi$ directly. Their
convergence is the content of CY-C, not a consequence of any
functoriality. At $d = 3$ for
$K3 \times E$ the routes form a pentagon of five intertwiners with
$R_2$ (Borcherds) as the source branch, and they stratify by the
generator-rank invariant $\rho^{R_i}$ (an algebraic invariant of
the presentation of each route's target) rather than by
$\kappa_{\mathrm{ch}}$. The compact total-space value is
$\kappa_{\mathrm{ch}}^{\mathrm{Hodge}}(K3\times E)=\chi(\cO_{K3\times E})=0$,
while the additive scalar
$\kappa_{\mathrm{ch}}^{\mathrm{Heis}}(K3\times E)
=\kappa_{\mathrm{ch}}^{\mathrm{Heis}}(K3)
+\kappa_{\mathrm{ch}}^{\mathrm{Heis}}(E)=2+1=3$
belongs to the relative Heisenberg/free-field route. The
manifold invariant $\kappa_{\mathrm{cat}} = \chi(\cO_{K3 \times E})
= 2 \cdot 0 = 0$ (K\"unneth with $\chi(\cO_E) = 0$) is likewise
route-independent. The genuinely route-distinguishing invariant is the
generator-rank $\rho^{R_i}$. The BKM-side object is the
Hall--Drinfeld double of the $K3\times E$ CoHA after the positive half,
negative half, Cartan, Hopf pairing, and completion are constructed; it
is not the Mukai self-mirror K3 Yangian branch.
\end{remark}

\begin{remark}[Generators and relations obstruction]
\label{rem:why-cy-c-hard}
The structural difficulty of CY-C is that constructing $G(\cC)$
from $\cC$ requires a generators-and-relations presentation
of the braided category attached to $\Phi(\cC)$: $\Rep^{\Etwo}(\Phi_2(\cC))$
at $d=2$ and $\cZ(\Rep^{\Eone}(\Phi_d(\cC)))$ at $d\geq3$. But $\Phi(\cC)$
is constructed via factorisation envelopes that are inherently
\emph{not} presented by generators and relations. The known
verified cases (K3 lattice, toric CY$_3$) all have
generators-and-relations presentations of $\Phi(\cC)$ available
through other means (lattice-VOA structure, Hall algebra
structure). The general CY$_3$ case lacks such presentations.
For $K3 \times E$, even when generators and
relations exist (for K3 $\times$ E), the natural invariant
distinguishing routes is the generator-lattice rank, not the
modular characteristic.
\end{remark}


\section{Quantum group A vs B/C/D distinctions}
\label{sec:qg-abcd-distinctions}

The four classical series exhibit type-dependent behaviour at
several points relevant to the chiral quantum group construction.
MC3 is proved on the evaluation-generated core for all simple
types via the five-family mechanism; this section records the
type-dependent distinctions in that setting.

\begin{proposition}[MC3 five-family mechanism: types A--D and exceptional]
\label{prop:mc3-five-family-types}
The Vol~II MC3 theorem on the evaluation-generated core is
extended from the type-A Baxter
constraint to a universal five-family mechanism covering all
simple types. The five families are:
\begin{enumerate}[label=(\roman*)]
 \item \emph{Type A} (asymptotic prefundamentals): generic
 evaluation modules at distinct spectral parameters give
 prefundamental modules whose limits as
 $u_i \to \infty$ generate the evaluation core;
 \item \emph{Types B, C, D} (reflection-equation Shapovalov):
 the orthogonal/symplectic evaluation modules satisfy a
 reflection equation in addition to the YBE, and the
 Shapovalov form of these modules generates the evaluation
 core;
 \item \emph{Uniform across all types} (Chari--Moura
 multiplicity-free $\ell$-weights): the
 $\ell$-weight space decomposition of evaluation modules is
 multiplicity-free at generic spectral parameters, ensuring
 the core is generated by the simple objects;
 \item \emph{Elliptic} (theta-divisor complement): for
 elliptic deformations of the evaluation modules, the
 $\Theta$-divisor complement provides the evaluation generator
 set;
 \item \emph{Super} (parity-balance for super-Yangian): the
 super-Yangian $Y(\frakg(m|n))$ has its evaluation core
 generated by parity-balanced modules
 (super-trace-zero condition).
\end{enumerate}
The Baxter constraint $b = a - 1/2$ is specific to
the asymptotic-prefundamental type-A presentation; it does not generalise
to types B/C/D, where the reflection-equation Shapovalov gives
an alternative generator set.
\end{proposition}

\begin{remark}[Type-A vs B/C/D in the chiral quantum group]
\label{rem:qg-type-a-vs-bcd}
The chiral quantum group $\Uq(\widehat{\frakg})^{\mathrm{ch}}$
has type-dependent structure:
\begin{enumerate}[label=(\alph*)]
 \item Type $A_{N-1}$: $\Uq(\widehat{\fsl}_N)^{\mathrm{ch}}$
 admits the standard RTT presentation with the rational
 $R$-matrix
 $R(u) = (u\,\mathrm{Id} - \Psi\,P)/(u - \Psi)$
 (additive Yang). The qdet is central and gives the
 $U(\widehat{\fgl}_N)^{\mathrm{ch}}$ extension.
 (The RTT sign convention, fixed by Vol~I, is
 $[t_{ij}, t_{kl}] = \Psi(\delta_{il} t_{kj} - \delta_{kj} t_{il})$
 in the additive presentation; the opposite sign appears in Molev's
 multiplicative $1 - P/u$ presentation.)
 \item Types $B_n, C_n, D_n$: $\Uq(\widehat{\frakg})^{\mathrm{ch}}$
 requires the orthogonal/symplectic $R$-matrix and
 satisfies an additional reflection equation. The Drinfeld
 second presentation differs from type A by extra
 generator-level conditions.
 \item Exceptional $G_2, F_4, E_6, E_7, E_8$: the chiral
 quantum group exists abstractly via the universal Yangian
 of Drinfeld but lacks a closed-form RTT presentation; the
 evaluation-generated core construction
 (Proposition~\ref{prop:mc3-five-family-types}) provides
 access via Chari--Moura $\ell$-weights.
\end{enumerate}
\end{remark}

\begin{remark}[Super-Yangian scope]
\label{rem:super-yangian-scope}
The super-Yangian $Y(\frakg(m|n))^{\mathrm{ch}}$ is
\textsc{conjectural} at the chiral level:
\begin{enumerate}[label=(\alph*)]
 \item The Drinfeld--Jimbo presentation of the classical
 super-Yangian $Y_\hbar(\frakg(m|n))$ is established
 (Khoroshkin, Lukierski, Tsuboi);
 \item The super-trace-zero condition gives a generalised
 RLL relation with the super-$R$-matrix replacing the bosonic
 one;
 \item The chiral globalisation $Y(\frakg(m|n))^{\mathrm{ch}}$
 on a curve $X$ is conjectured but not constructed; the
 main obstruction is that $\Sigma_n$-coinvariants do not
 commute with super-symmetrisation in the way required for
 the factorisation envelope;
 \item The super-complementarity identity
 $\kappa_{\mathrm{ch}} + \kappa_{\mathrm{ch}}^! = \max(m, n)$ is verified per-case for
 small ranks but lacks a canonical pairing
 (super-trace vs Berezinian; Vol~II open frontier F4).
\end{enumerate}
\end{remark}


\section{Magnificent Four at \texorpdfstring{$d = 4$}{d=4}: the scalar trace of \texorpdfstring{$\PhiFA_4(\Perf(\C^4))$}{Phi-FA-4(Perf(C4))}}
\label{sec:magnificent-four-d4}

Conjecture~\ref{conj:cy-c-constructed-loci}~(vii) leaves the quantum-group
target at $d \geq 4$ unconstructed; in general the Drinfeld-centre
category of a Stage-2 $\Eone$ specialisation need not be of
quantum-group type. The flat affine slice $\cC = \Perf(\C^4)$ admits a
sharp chain-level statement about $\PhiFA_4(\cC)$ itself: the
Stage-$1$ object is an $E_4$ holomorphic factorisation algebra, and the
Nekrasov--Piazzalunga--Kool--Rennemo Magnificent Four formula is its
toric Hilbert-scheme scalar trace.

\begin{theorem}[\texorpdfstring{$E_4$}{E4}-factorisation structure on \texorpdfstring{$\PhiFA_4(\Perf(\C^4))$}{Phi-FA-4(Perf(C4))}]
\label{thm:mag-four-e4-structure-qgr}
The infinity-category $\Perf(\C^4)$ of perfect complexes on affine
Calabi--Yau $4$-space carries the canonical CY$_4$ structure given by
the holomorphic volume form
$\Omega_{\C^4} = dz_1 \wedge dz_2 \wedge dz_3 \wedge dz_4$. The
Stage-1 factorisation envelope
$\PhiFA_4(\Perf(\C^4))$ of Chapter~\ref{ch:e1-chiral} is a factorisation
algebra on $\R^8 \simeq \C^4$ (Costello--Gwilliam 2021 Vol.~II
\S 4.10) carrying a canonical $E_4$ holomorphic factorisation structure.
The PTVV form on $\mathcal M(\Perf(\C^4))$ has shift $-2$, so functions
on the moduli stack carry the non-degenerate shifted Poisson bracket of
cohomological degree $2$.
The equivariant Hilbert-scheme trace is a map
\[
 \Tr^{M4}_{T_{\mathrm{CY}}}\colon
 K^{T_{\mathrm{CY}}}\!\left(\bigsqcup_{n\geq0}\Hilb^n(\C^4)\right)
 \longrightarrow \Q(t_1,t_2,t_3,t_4)\llbracket q\rrbracket ,
\]
and is not an operadic structure on $\PhiFA_4(\Perf(\C^4))$.
\end{theorem}

\begin{proof}
Hochschild--Kostant--Rosenberg on smooth affine CY$_4$ identifies
$\HH^\bullet(\Perf(\C^4)) \simeq \mathrm{PV}^\bullet(\C^4)$ with
the Schouten--Nijenhuis calculus. Kontsevich--Tamarkin $E_4$-formality
gives a chain-level $E_4$-algebra structure, unique up to the
$\mathrm{GRT}_1(\Q)$-torsor (Willwacher 2010). Costello--Gwilliam
Vol.~II \S 4.10 packages this $E_4$-structure globally as a
factorisation algebra on $\R^8$. Pantev--To\"en--Vaqui\'e--Vezzosi give
	the shift $2-d=-2$ on the moduli of objects; the non-degenerate inverse
	is the shifted Poisson bracket of degree $2$ on functions. The Hilbert-scheme trace
is obtained only after passing to the toric DT$_4$ fixed-locus
realisation of the positive Hall shadow.
\end{proof}

\begin{definition}[Magnificent Four partition function]
\label{def:zm4-qgr}
Let $T^4 = (\C^\times)^4$ act diagonally on $\C^4$ with equivariant
parameters $\epsilon_i = \log t_i$ subject to the Calabi--Yau condition
$\sum_{i=1}^{4} \epsilon_i = 0$, cutting out the torus
$T_{\mathrm{CY}} = \{(t_1, t_2, t_3, t_4) \in T^4 : t_1 t_2 t_3 t_4 = 1\}$.
Let $p$ be the instanton variable and let $s$ be the $U(1\mid 1)$
mass--Coulomb ratio. For a monomial $x$ write
$[x]=x^{1/2}-x^{-1/2}$. The rank-one Magnificent Four partition function
of Nekrasov--Piazzalunga~\cite{NekrasovPiazzalunga2019} is the signed
$K$-theoretic DT$_4$ fixed-point sum
\[
 Z^{M4}_{1}(p;t_i,s)
 \;=\; \sum_{\pi} p^{|\pi|}\,(-1)^{\mu_\pi}\,[-\mathsf v_\pi],
\]
where $\pi$ ranges over solid partitions and
$(-1)^{\mu_\pi}[-\mathsf v_\pi]$ is the Oh--Thomas localisation weight
with the Nekrasov--Piazzalunga orientation sign. Kool--Rennemo~\cite{KoolRennemo2025}
prove the closed plethystic form
\[
 Z^{M4}_{1}(p;t_i,s)
 \;=\;
 \operatorname{PE}\!\left(
 \frac{[t_1t_2][t_1t_3][t_2t_3][s]}
 {[t_1][t_2][t_3][t_4][p\sqrt{s}][p/\sqrt{s}]}
 \right),
 \qquad t_1t_2t_3t_4=1.
\]
\end{definition}

\begin{theorem}[Generating-function match]
\label{thm:mag-four-partition-match-qgr}
Under the toric DT$_4$ trace comparison, the equivariant
$T_{\mathrm{CY}}$-$K$-theoretic trace of the positive Hilbert-scheme
shadow is the Magnificent Four partition function:
\[
 \sum_{n\geq0}p^n\,
 \Tr^{M4}_{T_{\mathrm{CY}}}\!\left([\Hilb^n(\C^4)]^{\mathrm{vir}}\right)
 \;=\; Z^{M4}_{1}(p;t_i,s)
 \in K^{T_{\mathrm{CY}}\times \C_s^\times}(\mathrm{pt})\llbracket p\rrbracket .
\]
The identity is an exact formal identity in equivariant $K$-theory.
\end{theorem}

\begin{proof}
The $T_{\mathrm{CY}}$-fixed points of $\Hilb^n(\C^4)$ are solid
partitions $\pi$ of size $n$. The DT$_4$ virtual tangent complex gives
the Oh--Thomas localisation class at $\pi$~\cite{OhThomas2023};
Nekrasov--Piazzalunga write this class as
$(-1)^{\mu_\pi}[-\mathsf v_\pi]$. Kool--Rennemo prove the
Nekrasov--Piazzalunga formula by refining the local classes to
factorisable sheaves and applying Okounkov rigidity~\cite{KoolRennemo2025}. The Hall trace of
the positive Hilbert-scheme shadow is precisely the sum of these
localised classes.
\end{proof}

\begin{example}[Fixed-point support]
\label{ex:Z-m4-expansion}
The unweighted support of the trace is the solid-partition set:
\[
 \sum_{n\geq0}|\Hilb^n(\C^4)^{T_{\mathrm{CY}}}|p^n
 = 1 + p + 4p^2 + 10p^3 + 26p^4 + 59p^5 + \cdots.
\]
The partition function itself weights this support by the DT$_4$
orientation sign and the square-root Euler class.
\end{example}

\begin{proposition}[CY-D at \texorpdfstring{$d = 4$}{d=4} on compact CY$_4$]
\label{prop:cy-d-d4-compact-qgr}
For a compact simply-connected Calabi--Yau $4$-fold $X$ with
Hodge pattern $h^{0,0}(X) = h^{0,4}(X) = 1$ and
$h^{0,1}(X) = h^{0,2}(X) = h^{0,3}(X) = 0$,
\[
 \kappa_{\mathrm{ch}}^{\mathrm{Hodge}}(\PhiFA_4(\Perf X)) \;=\; \chi(\cO_X) \;=\; 2.
\]
In particular, CY-D at $d = 4$ on simply-connected compact CY$_4$ is a
theorem; it is not conjectural in this case.
\end{proposition}

\begin{proof}
For compact $X$, the Hodge supertrace of $\PhiFA_4(\Perf X)$ equals
$\chi(\cO_X)$ by Theorem~\ref{thm:cy-modular-characteristic}.
Hirzebruch--Riemann--Roch gives
$\chi(\cO_X) = (3 c_2^2 - c_4)/720 = 2$ for simply-connected
CY$_4$.
\end{proof}

\begin{remark}[Four-\texorpdfstring{$\kappa_\bullet$}{kappa-.} values on the flat slice and the sextic]
\label{rem:four-kappa-d4-split}
On the affine slice $\cC = \Perf(\C^4)$ the four invariants split as
\[
 \kappa_{\mathrm{cat}}(\C^4) = \chi(\cO_{\C^4}) = 1, \qquad
 \kappa_{\mathrm{ch}}^{\mathrm{Hodge}}(\PhiFA_4(\Perf \C^4)) = 0,
\]
with $\kappa_{\mathrm{BKM}}, \kappa_{\mathrm{fiber}}$ not applicable
(the flat affine slice carries no BKM). The discrepancy
$1 \neq 0$ is the Berezin-integration volume factor: the categorical
Euler includes the unit section while the Hodge supertrace subtracts
the volume normalisation. On the compact sextic $X_6 \subset \mathbb{P}^5$
(simply-connected CY$_4$),
$\kappa_{\mathrm{cat}}(X_6) = \kappa_{\mathrm{ch}}^{\mathrm{Hodge}}(\PhiFA_4(\Perf X_6)) = 2$
by Proposition~\ref{prop:cy-d-d4-compact-qgr} with
$\kappa_{\mathrm{BKM}}$ open because no $d = 4$ BKM has been
constructed, and the middle Hodge number is \(h^{2,2}(X_6)=1752\). The full fourth Betti number is
\[
b_4(X_6)=h^{4,0}+h^{3,1}+h^{2,2}+h^{1,3}+h^{0,4}
=1+426+1752+426+1=2606.
\]
The primitive \((2,2)\)-summand has dimension \(1751\), after removing the ambient class \(H^2\).
\end{remark}

\begin{proposition}[$S_4$-symmetry at \texorpdfstring{$d = 4$}{d=4}]
\label{prop:s4-symmetry-d4-qgr}
$Z^{M4}_{1}(p;t_i,s)$ is invariant under permutation of
$(t_1,t_2,t_3,t_4)$ subject to $t_1t_2t_3t_4=1$. The action lifts to the
orientation groupoid of $\PhiFA_4(\Perf(\C^4))$ and preserves the scalar
trace $\Tr^{M4}_{T_{\mathrm{CY}}}$. This is the $d = 4$ analogue of the
Miki $\Z/3$-triality at $d = 3$ on
$\mathrm{CoHA}(\C^3) = Y^+(\widehat{\fgl}_1)$ (Miki 2007;
Feigin--Jimbo--Miwa--Mukhin 2016).
\end{proposition}

\begin{proof}
The displayed plethystic numerator chooses the polarisation
$\{t_1t_2,t_1t_3,t_2t_3\}$ and is not, by itself, an $S_4$-symmetric
function. A permutation of the four coordinates changes this
polarisation and also changes the DT$_4$ orientation sign. Monavari's
canonical-vertex theorem~\cite{Monavari2022}, used in the
Kool--Rennemo proof, states that
the product $(-1)^{\mu_\pi}[-\mathsf v_\pi]$ is invariant under this
simultaneous change. Since the coordinate permutation sends
$\Omega_{\C^4}$ to $\mathrm{sgn}(\sigma)\Omega_{\C^4}$, the lift acts on
the orientation line rather than on a fixed oriented chart. The local
weights and hence their Hall trace are therefore invariant.
\end{proof}

\begin{conjecture}[Maulik--Okounkov envelope at \texorpdfstring{$d = 4$}{d=4}]
\label{conj:mo-envelope-d4-qgr}
There exists a $d = 4$ analogue of the Maulik--Okounkov stable envelope
(\texttt{arXiv:1211.1287}) on the CY$_4$ Hilbert schemes
$\{\Hilb^n(\C^4)\}_{n \geq 0}$, constructed from the
$T_{\mathrm{CY}}$-equivariant $K$-theory via the Nekrasov--Piazzalunga
sign rule, and satisfying:
\begin{enumerate}[label=(\roman*)]
 \item a chamber structure indexed by orderings of
 $(\epsilon_1, \epsilon_2, \epsilon_3, \epsilon_4)$ on the CY$_4$ cone
 $\sum \epsilon_i = 0$;
 \item Yang--Baxter equation for the resulting $R$-matrix
 $R^{M4}(u, v)$ on
 $\bigoplus_n K^{T_{\mathrm{CY}}}(\Hilb^n(\C^4))$;
 \item generating-function compatibility
 $\sum_n q^n \Tr_{K^{T_{\mathrm{CY}}}(\Hilb^n)}\! R^{M4}(u, v)
 = Z^{M4}(q) \cdot f(u, v)$ for an explicit rational prefactor.
\end{enumerate}
This is the $d = 4$ incarnation of the universal positive-geometry
grammar $Y^+(X) = H^\bullet_{\mathrm{eq}}(\cM^+_{\mathrm{eff}}(X), \phi_W)$
with $X = \C^4$ and equivariance stratum (i)~toric $T^4$.
\end{conjecture}

\begin{remark}[Three Stage-2 specialisations at \texorpdfstring{$d = 4$}{d=4}]
\label{rem:spch-d4}
The two-stage factorisation
$\Phi_d = \SpCh_{\Sigma_{d-1}, C} \circ \PhiFA_d$ at $d = 4$ needs a
three-cycle $\Sigma_3$ and a curve $C$. The natural specialisations are
(T1)~torus $\Sigma_3 = T^3$ with $C = S^1$, (T2)~sphere $\Sigma_3 = S^3$
with $C$ a great circle, and (T3)~Hopf fibration
$\Sigma_3 = S^3 \to S^2$ with $C$ a Hopf fibre. The three produce
distinct chiral algebras on $C$ from a common Stage-1 source.
$\PhiFA_4$ is one functor; the Stage-2 stratum produces a
family of specialisations, not a family of $\Phi_4$-applications.
Mutual compatibility of the three specialisations is open (the
$d = 4$ analogue of the pentagon at $d = 3$).
\end{remark}

\begin{remark}[No quantum-group target at \texorpdfstring{$d = 4$}{d=4}]
\label{rem:upgrade-relative-cy-c-vii}
At $d \geq 4$ a quantum-group target $G(\cC)$ is undefined in general.
Theorem~\ref{thm:mag-four-e4-structure-qgr} does \emph{not} construct
$G(\cC)$ at $d = 4$; it constructs the Stage-1 factorisation envelope
$\PhiFA_4(\cC)$ at $\cC = \Perf(\C^4)$ as an $E_4$ holomorphic
factorisation algebra and identifies its toric scalar trace. The
quantum-group target $G(\Perf(\C^4))$ remains unconstructed. The
Drinfeld centre
$\cZ(\Rep^{\Eone}(\SpCh_{\Sigma_3,C}(\PhiFA_4(\Perf(\C^4)))))$, after
a Stage-2 specialisation $(\Sigma_3,C)$, would carry the
$\Etwo$-braided structure eligible for quantum-group identification
at the representation-category level; it remains conjectural.
\end{remark}


\section{Cross-volume \texorpdfstring{$\kappa_\bullet$}{kappa-.}-stratification}
\label{sec:cross-volume-kappa-strat}

The $\kappa_\bullet$-stratification by CY dimension (Vol~III
\cite{VolI}) restricts to the
quantum-group inputs of this chapter as follows. The chiral
Hodge-supertrace invariant $\kappa_{\mathrm{ch}}^{\mathrm{Hodge}}$
is K\"unneth-multiplicative on compact products and coincides with
$\chi(\cO_X)$ only for $d = 2$ with $h^{1,0}(X) = 0$; the
Heisenberg specialisation $\kappa_{\mathrm{ch}}^{\mathrm{Heis}}$ is
a separate generator-rank scalar. At odd $d$ one has
$\chi(\cO_X)=0$ by Serre parity.

\begin{proposition}[\texorpdfstring{$\kappa_{\mathrm{ch}}$}{kappa-ch} on QG inputs]
\label{prop:kappa-strat-qg-inputs}
For the standard quantum-group inputs to $\Phi$:
\begin{enumerate}[label=(\roman*)]
 \item $d = 1$, elliptic curve $E$: the compact Hodge supertrace is
 $\kappa_{\mathrm{ch}}^{\mathrm{Hodge}}(\Phi(E))=\chi(\cO_E)=0$,
 while the Heisenberg specialisation has
 $\kappa_{\mathrm{ch}}^{\mathrm{Heis}}(E)=1$; the QG output is a rank-$2$
 lattice VOA on the Heisenberg route;
 \item $d = 2$, K3: $\kappa_{\mathrm{ch}} = \chi(\cO_{K3}) = 2$
 (since $h^{1,0} = 0$); the QG output is the Mukai-lattice
 Heisenberg sector of Proposition~\ref{prop:kl-k3-lattice};
 \item $d = 2$, abelian surface $A$:
 $\kappa_{\mathrm{ch}}^{\mathrm{Hodge}}(\Phi(A))=\chi(\cO_A)=0$,
 while $\kappa_{\mathrm{ch}}^{\mathrm{Heis}}(A)=2$ on the
 Heisenberg route; the $d = 2$ identification
 $\kappa_{\mathrm{ch}}=\chi(\cO_X)$ requires $h^{1,0}=0$ and
 therefore does not apply to abelian surfaces;
 \item $d = 3$, $\C^3$: the Schiffmann--Vasserot cohomological
 Hall algebra $\mathrm{CoHA}(\C^3)$ is an associative algebra,
 isomorphic to the positive half $Y^+(\widehat{\fgl}_1)$ of the
 affine Yangian; $\mathrm{CoHA}$ is not itself a chiral algebra,
 and an $\Eone$-chiral algebra $\Phi(\C^3)$ whose Drinfeld centre
 satisfies
 $\cZ(\Rep^{\Eone}(\Phi(\C^3))) \simeq \Rep_q(\widehat{\fgl}_1)$
 would give $\kappa_{\mathrm{ch}} = c = 1$ for a Heisenberg
 realization, conditional on CY-C at noncompact CY$_3$;
 \item $d = 3$, conifold: the resolved conifold is not a local
 surface; direct McKay computation gives
 $\kappa_{\mathrm{ch}}(\cA_{\mathrm{conifold}})=1$. Its CoHA is a
 positive Hall half in the toric chamber; the full
 $Y(\widehat{\fsl}_2)$ requires the completed Drinfeld double;
 \item $d = 3$, compact quintic: $\chi(\cO_{\text{quintic}}) = 0$
 by Serre parity, so the $d = 2$ identification
 $\kappa_{\mathrm{ch}} = \chi(\cO_X)$ cannot survive; the QG
 output is conjectural (CY-C open at general compact CY$_3$);
 \item $d = 3$, $K3 \times E$:
 the compact Hodge supertrace is
 $\kappa_{\mathrm{ch}}^{\mathrm{Hodge}}(K3 \times E)=0$, while the
 Heisenberg/Stage-$2$ specialisation has
 $\kappa_{\mathrm{ch}}^{\mathrm{Heis}}(K3 \times E)
 = \kappa_{\mathrm{ch}}^{\mathrm{Heis}}(K3)
 + \kappa_{\mathrm{ch}}^{\mathrm{Heis}}(E) = 2 + 1 = 3$;
 the Hall--Drinfeld/Borcherds branch attached to
 $\mathfrak g_{\Delta_5}$ has
 $\kappa_{\mathrm{BKM}}(\Delta_5) = \left.c_N(0)/2\right|_{N=1}=5$ by
 Gritsenko's denominator theorem
 (Theorem~\ref{thm:borcherds-weight-kappa-BKM-universal}).
\end{enumerate}
\end{proposition}

\begin{remark}[\texorpdfstring{$\kappa_\bullet$}{kappa-.}-values for \texorpdfstring{$K3 \times E$}{K3 x E}]
\label{rem:three-kappa-k3xe}
The four invariants separate cleanly. The manifold invariant is
$\kappa_{\mathrm{cat}}(K3 \times E) = \chi(\cO_{K3 \times E})
= \chi(\cO_{K3}) \cdot \chi(\cO_E) = 2 \cdot 0 = 0$ by K\"unneth,
vanishing because $\chi(\cO_E) = 0$ at odd $d = 1$. The
algebraization invariants come from two different constructions on
the same manifold: the Heisenberg-level chiral algebra gives
$\kappa_{\mathrm{ch}}^{\mathrm{Heis}} = 3$ by additivity across the
product, while the Hall--Drinfeld/Borcherds branch gives
$\kappa_{\mathrm{BKM}}(\Delta_5) = \left.c_N(0)/2\right|_{N=1} = 10/2 = 5$ from
the primitive Gritsenko form $\Delta_5$ (whose primitive Igusa square
is $\Phi_{10}^{\mathrm{un}}=\Delta_5^2$), by
Theorem~\ref{thm:borcherds-weight-kappa-BKM-universal}. The fibre invariant is the
Mukai-lattice rank $\kappa_{\mathrm{fiber}}=24$. The BKM-side object is
the Hall--Drinfeld double of $\mathrm{CoHA}(K3\times E)$, not the K3
Yangian; the K3 Yangian is the separate Mukai self-mirror branch. The
tempting additive decomposition
$\kappa_{\mathrm{BKM}} \stackrel{?}{=} \kappa_{\mathrm{ch}}^{\mathrm{Hodge}} + \chi(\cO_{\text{fiber}})
= 3 + 2 = 5$ is convention-mixing: the $3$ is the Heisenberg
specialisation, not the compact total-space supertrace
$\kappa_{\mathrm{ch}}^{\mathrm{Hodge}}(K3 \times E)=0$. Under the compact convention
the additive formula already fails at $N=1$, since
$5 \neq 0 + \chi(\cO_E) = 0$; across the CHL $\Z/N\Z$ orbifolds it
continues to fail because the Borcherds constant changes to
$c_N(0)/2$. The universal formula is
$\kappa_{\mathrm{BKM}}(\Phi_N)=c_N(0)/2$. The invariants are
controlled by four different functorialities: compact Hodge supertrace,
manifold Euler characteristic, BKM denominator, and fibre lattice. The arithmetic
identity $5 = 3 + 2$ at $N = 1$ is not a structural decomposition.
\end{remark}


\section{Chiral QG equivalence on the Koszul locus}
\label{sec:outlook-chiral-qg-koszul}

Which $\Eone$-chiral algebras admit a bar--cobar calculus fine
enough to compare the three quantum-group structures below? The
answer is a locus cut out by the existence of an ordered chiral
Koszul datum.

\begin{definition}[Ordered Koszul locus;
\ClaimStatusDefinitional]
\label{def:kosz-ord-vol3}
Let $X$ be a smooth algebraic curve. An \emph{ordered chiral
Koszul datum} for an augmented $\Eone$-chiral algebra $A$ on
$\Ran(X)$ (factorisation structure over the ordered configuration
spaces $\Conf^{\mathrm{ord}}_n(X)$) is a triple
$(C, \tau, F_\bullet)$:
\begin{enumerate}[label=(\roman*)]
 \item a complete conilpotent coalgebra $C$ in the ordered
 factorisation category on $\Ran(X)$;
 \item a twisting morphism $\tau \colon C \to A$ satisfying the
 Maurer--Cartan equation in the ordered convolution algebra;
 \item an exhaustive, separated filtration $F_\bullet$ on $C$
 compatible with $\tau$, each of whose conformal-weight/bar-length
 pieces is finite-dimensional over the ground field or locally
 free of finite rank over $\cO_X$, with transition maps
 satisfying the Mittag--Leffler condition,
\end{enumerate}
such that the coalgebra unit $C \to B^{\mathrm{ord}}_X(A)$ into
the ordered bar coalgebra
$B^{\mathrm{ord}}_X(A) = T^c(s^{-1}\bar{A})$ and the cobar counit
$\Omega^{\mathrm{ord}}_X(C) \to A$ are filtered
quasi-isomorphisms. The \emph{ordered Koszul locus}
\[
\mathrm{Kosz}^{\mathrm{ord}}(X)
\;\subset\;
\Eone\text{-}\mathrm{ChirAlg}(X)
\]
is the full subcategory of augmented $\Eone$-chiral algebras on
$\Ran(X)$ admitting an ordered chiral Koszul datum. Membership is
the existence of such a datum; the comparison maps in the
statements below are attached to a chosen datum. This is the
ordered ($\Eone$) counterpart of the symmetric chiral Koszul
pairs of Vol~I \cite{VolI}, restated here self-containedly.
\end{definition}

\begin{theorem}[Chiral QG determinations on coproduct-equipped
input; \ClaimStatusProvedElsewhere]
\label{thm:chiral-qg-equiv-datum-vol3}
Let $A \in \mathrm{Kosz}^{\mathrm{ord}}(X)$ underlie a chiral
Yangian datum $(A, S(z), \Delta^{\mathrm{ch}}, \{m_k\})$ of
Definition~\ref{def:chiral-yangian-datum-vol3}, so that the
coproduct is supplied as part of the datum. Then, on the
evaluation-generated core and for all simple types, the
$R$-matrix $r(z) \in \End(A \otimes A)((z))$ and the chain-level
$\Ainf$ operations $\{m_k\}_{k \geq 1}$ determine each other up
to chiral quasi-isomorphism, and each is compatible with the
supplied coproduct: this is the Vol~I content of MC3, proved via
the five-family mechanism \cite{VolI}. For Yangian input
$Y(\frakg)^{\mathrm{ch}}$, the compatibility of the pair
$(r(z), \{m_k\})$ with the supplied coproduct follows from
Drinfeld's second presentation together with Etingof--Kazhdan
flatness \cite{VolI}. The scope is exactly this: the coproduct is
input, never output.
\end{theorem}

\begin{conjecture}[Chiral QG equivalence on the ordered Koszul
locus; \ClaimStatusConjectured]
\label{thm:chiral-qg-equiv-vol3}
On the ordered Koszul locus $\mathrm{Kosz}^{\mathrm{ord}}(X)$ of
Definition~\ref{def:kosz-ord-vol3}, the three structures
\begin{enumerate}[label=(\arabic*)]
 \item $R$-matrix: $r(z) \in \End(A \otimes A)((z))$;
 \item $\Ainf$ structure: chain-level operations
 $\{m_k\}_{k \geq 1}$ on $A$;
 \item Ordered coproduct: $\Delta^{\mathrm{ord}} \colon A \to
 A \otimes A$ at depth $z = z_1 - z_2$
\end{enumerate}
mutually determine each other up to chiral quasi-isomorphism.
The determinations within the pair (1)--(2), and their
compatibility with a coproduct supplied as input, are
Theorem~\ref{thm:chiral-qg-equiv-datum-vol3}. The open content is
the reconstruction of (3) from (1)--(2): the ordered coproduct
probes the depth-two collision residue $r(z_1 - z_2)$ at
$z_1 \neq z_2$, while the $\Eone$-factorisation data underlying
(1)--(2) probe only the depth-one residue in the limit
$z_1 \to z_2$. Whether the depth-two datum is determined by the
depth-one datum is precisely the open forgetful-arrow problem of
Remark~\ref{rem:yangian-forgetful-arrow}; no case of this
reconstruction is proved.
\end{conjecture}

\begin{remark}[Concrete instances in Vol III]
\label{rem:vol3-concrete-qg-instances}
The chiral QG equivalence has three concrete Vol~III instances:
\begin{enumerate}[label=(\alph*)]
 \item \emph{$\fsl_2$ Yangian}: explicit $\alpha, \beta, \gamma$
 order-$\hbar$ computation including the triangle coherence
 $\alpha \circ \gamma \circ \beta = \mathrm{id}$ (Vol~I
 \cite{VolI});
 \item \emph{Affine $V_k(\fsl_2)$}: KL equivalence at
 $q = e^{\pi i / (k+2)}$ (Theorem~\ref{thm:kl-equivalence})
 gives the chiral QG datum at integer level $k$;
 \item \emph{Toroidal Yangian on $K3 \times E$}: the elliptic
 deformation, conditional on chain-level class-M topologization
 (Vol~I open frontier F3) and the K3 elliptic fibration
 structure of Chapter~\ref{ch:k3-quantum-toroidal}.
\end{enumerate}
\end{remark}

\begin{remark}[Vol I comparison theorems]
\label{rem:vol3-comparison-theorems}
The chapter cross-refers to:
\begin{enumerate}[label=(\alph*)]
 \item Conjecture~\ref{thm:chiral-qg-equiv-vol3} (chiral QG
 equivalence on the ordered Koszul locus,
 Definition~\ref{def:kosz-ord-vol3}), whose
 coproduct-equipped sub-case is
 Theorem~\ref{thm:chiral-qg-equiv-datum-vol3};
 \item Vol~II \cite{VolII} (formal deformation
 families identification for simple $\frakg$, and
 chain-level identification at fixed KZ associator for
 simple / reductive / solvable $\frakg$; covers $\fgl_N$,
 semisimple, nilpotent Heisenberg as subcases);
 \item Vol~I \cite{VolI}
 (gl$_N$ chiral QG via Feigin--Frenkel screening, all
 $N \geq 2$ unconditional);
 \item Vol~II \cite{VolII}
 (unified chiral quantum group $Q_g^{k,f,\mu}$ covering
 nine specialisation fibres);
\item the universal trace identity
$K(A) = -c_{\mathrm{ghost}}(\mathrm{BRST}(A))$ on the Vol~I
 side bridging to $\kappa_{\mathrm{BKM}}(\Phi_N) = c_N(0)/2$
 on the Vol~III side via $\Phi$ for K3-fibered Class A.
\end{enumerate}
The Vol~III incarnation of the chiral QG equivalence is the
chapter's own
Conjecture~\ref{conj:cy-c-constructed-loci} restricted to verified
cases (i)--(iv). Cases (v)--(vii) remain conditional on their comparison maps.
\end{remark}

\section{\texorpdfstring{$K3 \times E$}{K3 x E} Hall--Drinfeld double representation theory}
\label{sec:qgreps-supplement}

\begin{conjecture}[\texorpdfstring{$K3 \times E$}{K3 x E} Hall--Drinfeld double representation theory]
\label{rem:qgreps-hall-drinfeld-reps}
Assume the oriented critical CoHA positive half
$Y^+_{\mathrm{Hall}}(K3\times E)$ is constructed with its negative
half, Cartan part, nondegenerate Hall--Serre Hopf pairing, completion,
and compatibility with the Drinfeld-centre half-braiding of
$\mathbf{H}_{\Delta_5}$. Then the completed Hall--Drinfeld double
\[
 \mathcal{D}_\hbar\bigl(Y^+_{\mathrm{Hall}}(K3\times E)\bigr)
\]
admits evaluation modules $V_u$ on which the Siegel-elliptic dynamical
$R$-matrix $R_{\mathrm{Sieg,dyn}}(u-v;\tau,\rho,z)$ acts as the
half-braiding $\sigma_{V_u}(V_v)$ in
$\mathcal{Z}(\mathrm{Rep}^{\Eone}(\mathbf{H}_{\Delta_5}))$. At the
abelian Lie/Hopf level one sees the $24$ Mukai-lattice Heisenberg
directions. The non-abelian BKM structure belongs to the vertex closure
of the Borcherds denominator package, not to the Mukai self-mirror K3
Yangian branch.
\end{conjecture}
